% * =========================================================================
% * 存在的困难与问题
% * =========================================================================
\section{存在的困难与问题}

本课题在推进过程中，主要面临以下困难与问题：

\subsection{存在问题与困难}

\begin{semiQuotList}
    \item \textbf{特定配置下的性能瓶颈}：在$(10,4)$等配置中，第四个校验块需要2个码字的移位偏移，无法在AVX2的256位寄存器约束下与其他校验块合并，导致编码和解码性能相对下降，这是当前实现尚未完全克服的瓶颈。
    \item \textbf{两块丢失场景的优化空间}：两块丢失时既无法使用单块丢失的直接异或快路径，又不具备高丢失数场景下异或占比高的特点，必须完整执行代入与消除，因此整体性能优势相对有限。
    \item \textbf{优化方法论的通用性}：当前优化方案主要针对Two-tone码的特定移位结构，如何将其推广至更广泛的移位异或码族并形成通用方法论仍需深入研究。
    \item \textbf{工程实现与系统集成的复杂性}：将算法以插件形式稳定集成进Ceph这类大型分布式系统，涉及接口适配、内存管理与正确性验证等诸多工程细节，整体调试与测试成本较高。
    \item \textbf{小规模覆盖写的区间映射难点}：第二个创新点首先要把对象偏移和写入长度准确映射到数据块、条带单元、Two-tone窗口和移位后的校验区间，避免跨边界和非对齐处理导致校验更新位置错误。
    \item \textbf{增量更新策略的不确定性}：增量更新需要额外读取旧数据和旧校验，完整重新编码则访问更连续、SIMD利用率可能更高；二者在不同更新大小和$(k,m)$配置下的切换条件仍需验证。
    \item \textbf{公平比较与瓶颈归因问题}：ISA-L Reed-Solomon本身也支持高效增量更新，需要区分性能差异来自有限域乘加、移位异或、内存访问，还是Ceph事务、网络和存储I/O等系统开销。
\end{semiQuotList}

\subsection{解决方案}

针对上述困难，拟采取以下解决方案：
\begin{semiQuotList}
    \item 针对寄存器约束导致的瓶颈，探索更灵活的码字分组与SIMD指令调度策略，或结合AVX-512等更宽寄存器进一步挖掘合并潜力。
    \item 针对两块丢失场景，分析其访问模式特征，尝试设计介于单块快路径与通用消除之间的中间优化策略。
    \item 通过抽象Two-tone优化技术中与具体偏移结构无关的共性部分，逐步构建面向移位异或码族的通用优化框架。
    \item 充分利用Ceph官方的测试工具与社区文档，结合单元测试与端到端测试保证集成的正确性与稳定性。
    \item 针对区间映射难点，设计统一的range mapping模块，将一次局部写转换为若干连续的规范化更新任务，分别处理边界区间和中间SIMD XOR主体区间。
    \item 针对更新策略选择，通过不同更新大小、条带大小和$(k,m)$配置下的延迟、访存量和CPU周期实验寻找阈值，形成随工作负载变化的增量更新/完整重新编码自适应选择策略。
    \item 针对公平比较与瓶颈归因，按纯计算层、RMW读写层和小规模系统层开展分层性能剖析，给出Two-tone/FastTT增量更新的适用边界与瓶颈来源。
\end{semiQuotList}
